% !TeX program = xelatex
\documentclass[a4paper,10pt]{article}

% ============================================================================
% សន្លឹកកិច្ចការសិស្ស ស-ខ្មៅ (Student Worksheet - Grade 8 3.1.2 PROP2)
% បរិបទ៖ គ្មានស្លាយ LCD និងគ្មានម៉ាស៊ីនគិតលេខ — សិស្សបំពេញផ្ទាល់លើសន្លឹកកិច្ចការ ២ ទំព័រ A4
% ============================================================================
\usepackage[a4paper, left=0.85cm, right=0.85cm, top=0.70cm, bottom=0.95cm, footskip=0.48cm]{geometry}
\usepackage{amsmath, amssymb}
\usepackage{unicode-math}
\usepackage{fontspec}
\usepackage{array, colortbl}
\usepackage[table]{xcolor}
\usepackage{enumitem}
\usepackage{pifont}
\usepackage{fancyhdr}
\usepackage{lastpage}
\usepackage{tikz}
\usepackage[most]{tcolorbox}
\usetikzlibrary{calc, arrows.meta, decorations.pathreplacing, positioning, babel}

\makeatletter
\let\@expl@@@mark@update@singlecol@structures@@\relax
\makeatother

\providecommand{\NGSFindFonts}{%
  \IfFileExists{C:/Windows/Fonts/cambria.ttc}{%
    \def\NGSSysPath{C:/Windows/Fonts/}%
    \def\NGSUserPath{C:/Users/KIT/AppData/Local/Microsoft/Windows/Fonts/}%
  }{%
    \IfFileExists{./fonts/KhmerOSsiemreap.ttf}{\def\NGSSysPath{./fonts/}\def\NGSUserPath{./fonts/}}{%
    \IfFileExists{../fonts/KhmerOSsiemreap.ttf}{\def\NGSSysPath{../fonts/}\def\NGSUserPath{../fonts/}}{%
    \IfFileExists{../../fonts/KhmerOSsiemreap.ttf}{\def\NGSSysPath{../../fonts/}\def\NGSUserPath{../../fonts/}}{%
    \IfFileExists{../../../fonts/KhmerOSsiemreap.ttf}{\def\NGSSysPath{../../../fonts/}\def\NGSUserPath{../../../fonts/}}{%
    \IfFileExists{../../../../fonts/KhmerOSsiemreap.ttf}{\def\NGSSysPath{../../../../fonts/}\def\NGSUserPath{../../../../fonts/}}{%
    \IfFileExists{../../../../../fonts/KhmerOSsiemreap.ttf}{\def\NGSSysPath{../../../../../fonts/}\def\NGSUserPath{../../../../../fonts/}}{%
    \IfFileExists{../../../../../../fonts/KhmerOSsiemreap.ttf}{\def\NGSSysPath{../../../../../../fonts/}\def\NGSUserPath{../../../../../../fonts/}}{%
      \def\NGSSysPath{../../../../../../../fonts/}\def\NGSUserPath{../../../../../../../fonts/}%
    }}}}}}}%
  }%
}
\NGSFindFonts

\setmathfont[Path={\NGSSysPath}, FontIndex=1, Scale=1.02]{cambria.ttc}
\defaultfontfeatures{Ligatures=TeX}

\setmainfont[
  Path={\NGSUserPath}, Extension=.ttf, Script=Khmer,
  BoldFont={KhmerOSsiemreap}, BoldFeatures={FakeBold=1.8},
  ItalicFont={KhmerOSsiemreap}, ItalicFeatures={FakeSlant=0.22}
]{KhmerOSsiemreap}

\newfontfamily\muol[Path={\NGSUserPath}, Extension=.ttf, Script=Khmer]{KhmerOSmuollight}
\newfontfamily\TimesEn[Path={\NGSSysPath}, Extension=.ttf, BoldFont={timesbd}, ItalicFont={timesi}]{times}

\XeTeXlinebreaklocale "km"
\XeTeXlinebreakskip = 0pt plus 1pt minus 0.1pt

\definecolor{panelgray}{gray}{0.93}
\definecolor{headgray}{gray}{0.85}

\pagestyle{fancy}
\fancyhf{}
\renewcommand{\headrulewidth}{0pt}
\fancyfoot[L]{\fontsize{7.5pt}{9.5pt}\selectfont អនុវិទ្យាល័យព្រះក្រសាំង \en{|} គណិតវិទ្យាថ្នាក់ទី៨ \en{|} មេរៀនទី៣៖ ១.២. លក្ខណៈនៃសមាមាត្រស្រប (គ្រូ៖ ហុក ប៊ុនថែន)}
\fancyfoot[R]{\fontsize{7.8pt}{9.5pt}\selectfont ទំព័រ \en{\thepage/\pageref{LastPage}}}

\setlength{\parindent}{0pt}
\setlength{\parskip}{0pt}
\setlength{\tabcolsep}{3.8pt}

\newcommand{\en}[1]{{\TimesEn #1}}
\newcommand{\enbf}[1]{{\TimesEn\bfseries #1}}
\newcommand{\unit}[1]{\,\text{\en{#1}}}
\newcommand{\cnum}[1]{\mbox{\fontsize{10pt}{12pt}\selectfont\ding{\numexpr181+#1\relax}}}
\newcommand{\chkbox}{\tikz[baseline=-0.12ex]\draw[black, line width=0.8pt] (0,0) rectangle (0.24cm,0.24cm);}
\newcommand{\blank}[1][1.5cm]{\makebox[#1]{\dotfill}}
\newcommand{\ansbox}[1][0.85cm]{\,\tikz[baseline=-0.5ex]\draw[black, line width=0.65pt, fill=white] (0,0) rectangle (#1, 0.42cm);\,}
\newcommand{\secbar}[2]{%
  \begin{tcolorbox}[enhanced, sharp corners, colback=headgray, colframe=black, boxrule=0.9pt, boxsep=1pt, top=2pt, bottom=2pt, left=5pt, right=5pt, before skip=4pt, after skip=3pt]
    {\muol\fontsize{9.2pt}{13pt}\selectfont #1}\hfill{\fontsize{8pt}{11pt}\selectfont\textbf{#2}}
  \end{tcolorbox}%
}

\newcolumntype{L}[1]{>{\raggedright\arraybackslash}p{#1}}
\newcolumntype{C}[1]{>{\centering\arraybackslash}p{#1}}

\begin{document}
\fontsize{9.6pt}{14.4pt}\selectfont

% ============================================================================
% ទំព័រទី ១៖ E1 + E2 (POE 1 & POE 2) + E3 សំយោគលក្ខណៈទាំងពីរ និងធនាគារហេតុផល
% ============================================================================
\begin{tcolorbox}[enhanced, sharp corners, colback=white, colframe=black, boxrule=1.0pt, boxsep=2pt, top=3pt, bottom=3pt, left=6pt, right=6pt]
  \begin{center}
    {\muol\fontsize{10.3pt}{15pt}\selectfont អនុវិទ្យាល័យព្រះក្រសាំង \enbf{--} សន្លឹកកិច្ចការសិស្ស \enbf{(}មេរៀនទី៣៖ ១\enbf{.}២\enbf{.} លក្ខណៈនៃសមាមាត្រស្រប\enbf{)}}\\[3pt]
    {\fontsize{8.8pt}{12pt}\selectfont
    ឈ្មោះសិស្ស / ដៃគូ៖ \blank[5.2cm] \hspace{0.25cm}
    ថ្នាក់ទី៨៖ \blank[1.1cm] \hspace{0.25cm}
    ក្រុមទី៖ \blank[1.0cm] \hspace{0.25cm}
    កាលបរិច្ឆេទ៖ ២៦/១១/២០២៦}
  \end{center}
\end{tcolorbox}

% --- E1 + E2 ---
\secbar{ផ្នែកទី១៖ បេសកកម្មទស្សន៍ទាយ \enbf{(E1)} និងការរុករកជាដៃគូ \enbf{(E2 -- POE 1 \& POE 2)}}{រយៈពេល ១៤ នាទី}

\begin{tcolorbox}[enhanced, sharp corners, colback=panelgray, colframe=black, boxrule=0.75pt, boxsep=1.5pt, top=2.5pt, bottom=2.5pt, left=5pt, right=5pt]
  \textbf{បេសកកម្ម \en{E1} (គូសធីក \chkbox\ មុនពិសោធន៍)៖}
  (១) ចតុកោណកែងសមាមាត្រ \en{A}$(9,6)$ និង \en{B}$(x,8)$ — តើបណ្ដោយ $x$ ស្មើ៖ \chkbox\ $11$ \quad \chkbox\ $12$\\
  (២) អង្ករ $320\unit{kg}$ ចែកតាមសមាជិក ៣ គ្រួសារ $(4,5,7)$ — តើគ្រួសារមាន $7$ នាក់ទទួលបាន៖ \chkbox\ $\approx 106.7\unit{kg}$ \quad \chkbox\ $140\unit{kg}$
\end{tcolorbox}

\vspace{3pt}
\textbf{ក. សកម្មភាពទី ១ (\en{POE 1})៖ ស្ពានតួចុង-តួមធ្យម លើចតុកោណកែងសមាមាត្រ \en{A, B, C} (សៀវភៅពុម្ព ទំ.២៨–២៩)}

{\renewcommand{\arraystretch}{1.48}
\begin{tabular}{|C{3.30cm}|L{5.65cm}|L{5.25cm}|C{3.40cm}|}
  \hline
  \rowcolor{headgray}
  \textbf{គូចតុកោណកែង} & \textbf{សមាមាត្រ $\dfrac{\text{ទទឹង}}{\text{បណ្ដោយ}}$ (ឯកតាត្រូវគ្នា)} & \textbf{ផលគុណតួចុង $=$ ផលគុណតួមធ្យម} & \textbf{តម្លៃមិនស្គាល់} \\
  \hline
  \en{A}$(9,6)$ និង $(6,4)$ & $\dfrac{6}{9} = \dfrac{4}{6}$ & $6\times 6 = 36$ \;និង\; $9\times 4 = 36$ & ស្មើគ្នាពិត (\ding{51}) \\
  \hline
  \en{A}$(9,6)$ និង \en{B}$(x,8)$ & $\dfrac{6}{9} = \dfrac{8}{x}$ & $6\cdot x = 9\times 8 = \blank[1.2cm]$ & $x = \blank[1.5cm]$ \\
  \hline
  \en{A}$(9,6)$ និង \en{C}$(15,y)$ & $\dfrac{6}{9} = \dfrac{y}{15}$ & $9\cdot y = 6\times 15 = \blank[1.2cm]$ & $y = \blank[1.5cm]$ \\
  \hline
  \en{B}$(12,8)$ និង \en{D}$(18,z)$ & $\dfrac{8}{12} = \dfrac{z}{18}$ & $12\cdot z = 8\times \blank[0.9cm] = \blank[1.1cm]$ & $z = \blank[1.5cm]$ \\
  \hline
\end{tabular}}

\vspace{3pt}
$\bullet$ \textbf{ពន្យល់ (\en{E})៖} ក្នុងសមាមាត្រ $\dfrac{y_1}{x_1} = \dfrac{y_2}{x_2}$ បើគុណអង្គទាំងពីរនឹង $x_1 x_2$ យើងទទួលបាន៖ $x_2 y_1 = \blank[2.0cm]$ ។

\vspace{4pt}
\textbf{ខ. សកម្មភាពទី ២ (\en{POE 2})៖ ពិសោធន៍ផលបូកផលធៀបស្មើគ្នា (សៀវភៅពុម្ព ទំ.២៩)}

{\renewcommand{\arraystretch}{1.48}
\begin{tabular}{|L{4.80cm}|C{2.45cm}|C{2.45cm}|C{2.45cm}|C{5.25cm}|}
  \hline
  \rowcolor{headgray}
  \textbf{បរិមាណ} & \textbf{គូទី ១} & \textbf{គូទី ២} & \textbf{គូទី ៣} & \textbf{ផលបូកសរុបទាំង ៣ គូ} \\
  \hline
  ចំនួនសមាជិក $x_i$ (នាក់) & $x_1 = 4$ & $x_2 = 5$ & $x_3 = 7$ & $x_1+x_2+x_3 = 4+5+7 = \blank[0.9cm]$ \\
  \hline
  ទម្ងន់អង្ករ $y_i$ ($\unit{kg}$) & $y_1 = 80$ & $y_2 = 100$ & $y_3 = 140$ & $y_1+y_2+y_3 = \blank[1.8cm]$ \\
  \hline
  ផលធៀប $\dfrac{y_i}{x_i}$ ($\unit{kg}$/នាក់) & $\dfrac{80}{4} = 20$ & $\dfrac{100}{5} = \blank[0.8cm]$ & $\dfrac{140}{7} = \blank[0.8cm]$ & $\dfrac{y_1+y_2+y_3}{x_1+x_2+x_3} = \dfrac{320}{\blank[0.8cm]} = \blank[0.9cm]$ \\
  \hline
\end{tabular}}

\vspace{3pt}
$\bullet$ \textbf{ពន្យល់ (\en{E})៖} ផលធៀបនៃផលបូកភាគយក ចែកនឹងផលបូកភាគបែង មានតម្លៃ \blank[1.8cm] នឹងផលធៀបដើមនីមួយៗជានិច្ច!

% --- E3 ---
\secbar{ផ្នែកទី២៖ សំយោគលក្ខណៈទាំងពីរ \enbf{(E3)} និងធនាគារហេតុផល \enbf{B\&W} \cnum{1}\enbf{--}\cnum{5}}{សៀវភៅពុម្ព ទំ.២៨–២៩}

\begin{tabular}{|L{9.20cm}|L{9.20cm}|}
  \hline
  \textbf{១. លក្ខណៈទាំងពីរនៃសមាមាត្រស្រប (ទំ.២៨–២៩)៖}\par\vspace{2pt}
  $\bullet$ \textbf{លក្ខណៈទី ១ (តួចុង-តួមធ្យម)៖} បើ $\dfrac{y_1}{x_1} = \dfrac{y_2}{x_2}$ នោះ៖\par\vspace{2pt}
  {\centering $x_2 \cdot y_1 = \ansbox[1.4cm]$ \quad (ផលគុណតួចុង $=$ ផលគុណតួមធ្យម)\par}\vspace{2pt}
  $\bullet$ \textbf{លក្ខណៈទី ២ (ផលបូកផលធៀបស្មើគ្នា)៖}\par\vspace{2pt}
  {\centering $\dfrac{y_1}{x_1} = \dfrac{y_2}{x_2} = \dfrac{y_3}{x_3} = \dfrac{y_1 + y_2 + y_3}{\blank[2.4cm]}$\par}
  &
  \textbf{២. ធនាគារហេតុផល \en{B\&W} (\cnum{1}--\cnum{5})៖}\par\vspace{2pt}
  \cnum{1} គុណខ្វែងតួចុង-តួមធ្យម៖ $\frac{a}{b}=\frac{c}{d} \iff ad = bc$ ។\par\vspace{2pt}
  \cnum{2} ឯកតាជួរដេក៖ ដាក់ឯកតាដូចគ្នាក្នុងជួរដេកតែមួយមុនគុណខ្វែង ។\par\vspace{2pt}
  \cnum{3} ផលបូកផលធៀបស្មើគ្នា៖ $\frac{x}{a}=\frac{y}{b}=\frac{z}{c}=\frac{x+y+z}{a+b+c}$ ។\par\vspace{2pt}
  \cnum{4} គំរូរបារភាគស្មើគ្នា៖ បូកចំនួនចំណែកសរុប រួចចែករកតម្លៃ $1$ ចំណែក ។\par\vspace{2pt}
  \cnum{5} ផ្ទៀងផ្ទាត់ផលបូក៖ $x+y+z$ ត្រូវស្មើបរិមាណសរុបដើមជានិច្ច! \\
  \hline
\end{tabular}

\vspace{3pt}
\begin{tcolorbox}[enhanced, sharp corners, colback=panelgray, colframe=black, boxrule=0.65pt, boxsep=1.5pt, top=2pt, bottom=2pt, left=5pt, right=5pt]
  \textbf{ប្រយ័ត្នអន្ទាក់៖} (១) ក្នុងសមាមាត្រ $\frac{\text{ថ្លៃ}}{\text{ចំនួនផ្លៃ}} = \frac{\text{ថ្លៃ}}{\text{ចំនួនផ្លៃ}}$ ហាមដាក់ឯកតាឆ្លាស់ជួរ! \quad (២) ពេលចែកបរិមាណតាមសមាមាត្រ ត្រូវបូកចំនួនចំណែកសរុបជាមុនសិន!
\end{tcolorbox}

\vspace{3pt}
\textbf{៣. លំហាត់ត្រួតពិនិត្យរហ័ស៖ គណនាតួមិនស្គាល់ និងបែងចែកបរិមាណតាមសមាមាត្រ}

\vspace{2pt}
{\renewcommand{\arraystretch}{1.64}
\begin{tabular}{|C{4.40cm}|C{4.40cm}|C{4.40cm}|C{4.40cm}|}
  \hline
  \rowcolor{panelgray}
  $\dfrac{3}{5} = \dfrac{x}{20} \Rightarrow x = \blank[1.0cm]$ & $\dfrac{12}{x} = \dfrac{4}{7} \Rightarrow x = \blank[1.0cm]$ & $\dfrac{x}{2}=\dfrac{y}{3}=\dfrac{25}{5} \Rightarrow x=\blank[0.8cm]$ & $\dfrac{x}{2}=\dfrac{y}{3}=5 \Rightarrow y=\blank[0.8cm]$ \\
  \hline
  $\dfrac{7}{10} = \dfrac{21}{x} \Rightarrow x = \blank[1.0cm]$ & $\dfrac{x}{15} = \dfrac{8}{12} \Rightarrow x = \blank[1.0cm]$ & ចែក $60$ តាម $1\!:\!2 \Rightarrow \blank[1.1cm]$ & ចែក $100$ តាម $2\!:\!3 \Rightarrow \blank[1.0cm]$ \\
  \hline
  $\dfrac{5}{8} = \dfrac{25}{x} \Rightarrow x = \blank[1.0cm]$ & $\dfrac{9}{x} = \dfrac{3}{11} \Rightarrow x = \blank[1.0cm]$ & $\dfrac{x}{3}=\dfrac{y}{5}=\dfrac{40}{8} \Rightarrow x=\blank[0.8cm]$ & $\dfrac{x}{3}=\dfrac{y}{5}=5 \Rightarrow y=\blank[0.8cm]$ \\
  \hline
  $\dfrac{x}{14} = \dfrac{6}{7} \Rightarrow x = \blank[1.0cm]$ & $\dfrac{18}{24} = \dfrac{x}{8} \Rightarrow x = \blank[1.0cm]$ & ចែក $80$ តាម $3\!:\!5 \Rightarrow \blank[1.1cm]$ & ចែក $90$ តាម $2\!:\!3\!:\!4 \Rightarrow \blank[0.9cm]$ \\
  \hline
  $\dfrac{4}{9} = \dfrac{16}{x} \Rightarrow x = \blank[1.0cm]$ & $\dfrac{15}{x} = \dfrac{5}{12} \Rightarrow x = \blank[1.0cm]$ & ចែក $120$ តាម $1\!:\!3 \Rightarrow \blank[1.0cm]$ & ចែក $150$ តាម $2\!:\!3 \Rightarrow \blank[1.0cm]$ \\
  \hline
\end{tabular}}

\newpage

% ============================================================================
% ទំព័រទី ២៖ E4 លំហាត់គំរូ + E5 ប្រតិបត្តិ + វិនិច្ឆ័យ A-D + Exit Check + កិច្ចការផ្ទះ
% ============================================================================
\secbar{ផ្នែកទី៣៖ ពង្រីកចំណេះដឹងតាមគំរូរបារភាគស្មើគ្នា \enbf{(E4 --} គំរូ ទំ\enbf{.}២៩\enbf{--}៣០\enbf{)}}{រយៈពេល ១០ នាទី}

\textbf{លំហាត់គំរូ (សៀវភៅពុម្ព ទំ.២៩–៣០)៖} កាកបាទក្រហមកម្ពុជាចែកអង្ករ $\mathbf{320\unit{kg}}$ ឱ្យប្រជាជន ៣ គ្រួសារតាមសមាមាត្រនៃចំនួនសមាជិក គឺគ្រួសារទី១ មាន $4$ នាក់, គ្រួសារទី២ មាន $5$ នាក់, និងគ្រួសារទី៣ មាន $7$ នាក់ ។ រកទម្ងន់អង្ករដែលគ្រួសារនីមួយៗទទួលបាន៖

{\renewcommand{\arraystretch}{1.48}
\begin{tabular}{|L{4.70cm}|L{8.90cm}|C{4.20cm}|}
  \hline
  \rowcolor{headgray}
  \textbf{ជំហានស៊ើបអង្កេត} & \textbf{ការសរសេរសមាមាត្រ និងគំរូរបារភាគស្មើគ្នា} & \textbf{លទ្ធផល និងហេតុផល} \\
  \hline
  \textbf{១. តាងអញ្ញាត និងផលបូកសរុប} & តាង $x, y, z$ ជាអង្ករគ្រួសារទាំង ៣៖ $x + y + z = \blank[1.4cm]\unit{kg}$ ; សមាជិកសរុប $4+5+7 = \blank[0.9cm]$ នាក់ & $16$ ចំណែកស្មើគ្នា (\cnum{4}) \\
  \hline
  \textbf{២. អនុវត្តលក្ខណៈទី ២} & $\dfrac{x}{4} = \dfrac{y}{5} = \dfrac{z}{7} = \dfrac{x+y+z}{4+5+7} = \dfrac{320}{\blank[0.9cm]} = \blank[1.1cm]\unit{kg}\text{/នាក់}$ & $1$ នាក់បាន $\blank[0.9cm]\unit{kg}$ (\cnum{3}) \\
  \hline
  \textbf{៣. គណនាចំណែកគ្រួសារនីមួយៗ} & $x = 4\times \blank[0.8cm] = \blank[1.0cm]$ ; $y = 5\times \blank[0.8cm] = \blank[1.0cm]$ ; $z = 7\times \blank[0.8cm] = \blank[1.0cm]$ & ផ្ទៀង៖ $\blank[0.8cm]+\blank[0.8cm]+\blank[0.8cm]=320$ (\cnum{5}) \\
  \hline
\end{tabular}}

\vspace{4pt}
$\bullet$ \textbf{សំណួរឆ្លុះបញ្ចាំង \en{(EIA Looking-Back)}៖} បើសិស្សម្នាក់យក $320 \div 3 \approx 106.7\unit{kg}$ តើគាត់ខុសត្រង់ហេតុផលលេខប៉ុន្មាន? \blank[3.2cm]

% --- E5 ---
\secbar{ផ្នែកទី៤៖ វាយតម្លៃលទ្ធផលសិក្សា \enbf{(E5 --} ប្រតិបត្តិ ១ និង ២ \enbf{+} វិនិច្ឆ័យ \enbf{A--D +} សំបុត្រចេញ\enbf{)}}{រយៈពេល ១០ នាទី}

\textbf{ក. លំហាត់ប្រតិបត្តិ (សៀវភៅពុម្ព ទំ.២៩ និង ទំ.៣០ លើ — ធ្វើជាដៃគូ រួចប្ដូរសន្លឹកផ្ទៀង \en{Pause-Review})៖}

{\renewcommand{\arraystretch}{1.48}
\begin{tabular}{|L{8.90cm}|L{8.90cm}|}
  \hline
  \textbf{១) ប្រតិបត្តិទី ១ (ទំ.២៩ — ផ្លែក្រូច)៖} ផ្លែក្រូច $9$ ផ្លៃ ថ្លៃ $1200\text{៛}$ ។\par\vspace{2pt}
  $\bullet$ \textbf{ក.} រកថ្លៃក្រូច $12$ ផ្លៃ៖ $\dfrac{1200}{9} = \dfrac{y}{12}$\par\vspace{2pt}
  $\implies 9y = 1200\times 12 = \blank[1.3cm] \implies y = \blank[1.4cm]\text{៛}$ (\cnum{1}, \cnum{2})\par\vspace{2pt}
  $\bullet$ \textbf{ខ.} បើមានប្រាក់ $32\,000\text{៛}$ ទិញក្រូចបានប៉ុន្មានផ្លៃ?\par\vspace{2pt}
  $\dfrac{1200}{9} = \dfrac{32\,000}{x} \implies x = \dfrac{9\times 32\,000}{1200} = \blank[1.4cm]$ ផ្លៃ
  &
  \textbf{២) ប្រតិបត្តិទី ២ (ទំ.៣០ — ចែកប្រាក់កូន ៣ នាក់)៖}\par\vspace{2pt}
  ចែកប្រាក់ $140\,000\text{៛}$ តាមសមាមាត្រអាយុកូន $7, 10$ និង $11$ ឆ្នាំ៖\par\vspace{2pt}
  $\bullet$ ផលបូកចំណែកអាយុ៖ $7 + 10 + 11 = \blank[1.0cm]$ ចំណែក\par\vspace{2pt}
  $\bullet$ $\dfrac{x}{7}=\dfrac{y}{10}=\dfrac{z}{11}=\dfrac{140\,000}{\blank[0.9cm]} = \blank[1.4cm]\text{៛}$\par\vspace{2pt}
  $\bullet$ $x = \blank[1.4cm]\text{៛}$ ; $y = \blank[1.4cm]\text{៛}$ ; $z = \blank[1.4cm]\text{៛}$ (\cnum{5}) \\
  \hline
\end{tabular}}

\vspace{5pt}
\textbf{ខ. គូសវង់ជុំវិញជម្រើសត្រឹមត្រូវតែ ១ គត់ ក្នុងចំណោម \en{A, B, C, D} រួចសរសេរលេខហេតុផល \cnum{1}--\cnum{5} នៃជម្រើសខុស៖}

{\renewcommand{\arraystretch}{1.36}
\begin{tabular}{|C{4.40cm}|C{4.40cm}|C{4.40cm}|C{4.40cm}|}
  \hline
  \cellcolor{panelgray}\textbf{ជម្រើស \en{A}} & \cellcolor{panelgray}\textbf{ជម្រើស \en{B}} & \cellcolor{panelgray}\textbf{ជម្រើស \en{C}} & \cellcolor{panelgray}\textbf{ជម្រើស \en{D}} \\
  \hline
  កូនអាយុ $11$ បាន $55\,000\text{៛}$ & $\frac{6}{9}=\frac{8}{x} \Rightarrow 8x=54$ & អង្ករ $320\unit{kg}$ ចែក $\frac{320}{3}\unit{kg}$ & ក្រូច៖ $\frac{9}{1200}=\frac{1600}{x}$ \\[3pt]
  \chkbox\ ត្រូវ \quad \chkbox\ ខុស (\blank[0.7cm]) & \chkbox\ ត្រូវ \quad \chkbox\ ខុស (\blank[0.7cm]) & \chkbox\ ត្រូវ \quad \chkbox\ ខុស (\blank[0.7cm]) & \chkbox\ ត្រូវ \quad \chkbox\ ខុស (\blank[0.7cm]) \\
  \hline
\end{tabular}}

\vspace{5pt}
\begin{tcolorbox}[enhanced, sharp corners, colback=white, colframe=black, boxrule=0.85pt, boxsep=2.5pt, top=3.5pt, bottom=3.5pt, left=5pt, right=5pt]
  {\muol\fontsize{8.8pt}{12.5pt}\selectfont សំបុត្រចេញរហ័ស \enbf{(Exit Check --} ៦០ វិនាទី\enbf{)}៖}\\[2pt]
  (១) រក $x$ បើ $\dfrac{4}{6} = \dfrac{10}{x} \implies x = \blank[1.3cm]$ \;\en{|}\;
  (២) ចែកប្រាក់ $90\,000\text{៛}$ តាមសមាមាត្រ $2:3:4 \implies$ បាន \blank[1.5cm], \blank[1.5cm], \blank[1.5cm]
\end{tcolorbox}

\vspace{5pt}
\begin{tcolorbox}[enhanced, sharp corners, colback=panelgray, colframe=black, boxrule=0.85pt, boxsep=2.5pt, top=3.5pt, bottom=3.5pt, left=5pt, right=5pt]
  {\muol\fontsize{8.8pt}{12.5pt}\selectfont កិច្ចការផ្ទះ \enbf{(}សៀវភៅពុម្ព ទំព័រ ៤០\enbf{--}៤១ លំហាត់ \enbf{\#1, \#2, \#4, \#5}ក\enbf{, \#11} និង \enbf{\#12)}៖}\\[2pt]
  $\bullet$ \textbf{លំហាត់ \#1, \#2, \#4, \#5ក (ទំ.៤០)៖} សម្រួលផលធៀប និងរកតួមិនស្គាល់តាមផលគុណតួចុង-តួមធ្យម ។\\
  $\bullet$ \textbf{លំហាត់ \#11 និង \#12 (ទំ.៤១)៖} ដោះស្រាយចំណោទបែងចែកសមាមាត្រ និងអានមុន \textbf{១.៣. ទំហំសមាមាត្រច្រាស (ទំ.៣០–៣១)} ។
\end{tcolorbox}

\end{document}
